\label{sec:robust_hazard}

Under term-level scoring, leading filings appear to belong to better
firms, and the appearance is consistent across three outcomes. Among
registered debut filings, the probability that the owner ever appears in
SEC EDGAR rises in term-\dkl{} from 0.29\% at the bottom quintile to
0.42\% at the top. On an SEC-matched firm-year panel, a firm's trailing
3-year mean term-\dkl{} predicts its gross margin at $+2.2$pp per
standard deviation ($t = 2.7$ on the locally rebuilt panel through 2019;
$+3.2$pp, $t = 5.3$, on the published panel through 2024), with
past-facing surprise entering positively and future-facing negatively.
And firm-year term-\dkl{} predicts excess stock returns of $+1.5$ to
$+5.0$pp per standard deviation at one-to-four-year horizons. (A larger
debut-only return figure, $+28$pp at 4 years, fails a sample-size
diagnostic and is not claimed.)

None of these signed results survives when the same text is scored on
themes instead of words. On identical samples, the debut listing
gradient in signed theme-\dkl{} is flat to declining. The margin
coefficient reverses sign --- $-1.3$pp per standard deviation
($t = -1.8$) at $T = 50$, $-2.2$pp ($t = -3.1$) at $T = 200$, and
$-0.8$pp ($t = -1.1$) at $T = 500$, against the term-level $+2.3$pp ---
and the past-/future-facing decomposition flips sign and stays flipped.
The term-level listing gradient separately dissolves under
document-length controls: within bands of equal length it is U-shaped
rather than rising, and the monotone rise is contributed by composition
across lengths, because listed firms' counsel write longer descriptions.
The mark-level results, by contrast, are unaffected --- they hold under
both scorings at every resolution.

One rescue for the signed signal must be ruled out. The theme model's
vocabulary is built from a sample, so rare and emergent terms are
invisible to it; if the firm signal lived exactly in those terms, theme
scoring would be erasing a real signal rather than an artifact. To test
this, each filing's share of vocabulary invisible to the theme model is
computed --- 63.8\% of the software class's term \emph{types}, though
only about 9\% of its token \emph{mass}. That share is itself a strong
firm marker: P(EDGAR) runs from 2.9\% to 5.5\% across its quintiles, and
survival at the five-year proof improves by $+5.7$pp. But \emph{within}
strata of equal invisible share, the signed \dkl{} gradient on the firm
outcome is flat to negative, while the mark-level penalty holds in both
strata ($-6.9$ and $-8.3$pp). The firm-level signal that term scoring
finds is unsigned specificity: serious firms describe specific products
with specific, therefore rare, words, and rare-term mass inflates
term-\dkl{} mechanically. Note that raising the theme count does not
test this --- a crypto/blockchain theme exists at $T = 200$ but not at
$T = 50$ or $T = 500$, because raising $T$ re-partitions the same fixed
vocabulary rather than extending it --- so resolution sweeps test
granularity, not coverage, and the invisible-share split is the test
that settles the question.

The registration filter does not explain the divergence either. Dropping
the registration conditioning leaves the term-level listing gradient
essentially unchanged (0.26\% to 0.38\% across quintiles against 0.29\%
to 0.42\% conditional), and an examiner channel touching 4.3\% of failed
applications cannot generate several-point quintile gaps.

The hazard, stated for reuse: term-resolved text-novelty measures can
manufacture firm-performance correlations out of drafting style and
lexical specificity, and a distribution-based rescoring --- one model
fit --- is a cheap, decisive diagnostic.
